\documentclass[../../main.tex]{subfiles}
\graphicspath{{../../image/}{../../figure/}}
\begin{document}

\section{非守恒激波}\label{section:ch08-非守恒激波}

在连续介质物理中，人们偶尔会遇到非守恒形式的拟线性双曲系统
\begin{align}
\partial_t U(x,t)+A(U(x,t))\partial_x U(x,t)=0
\label{align:ch08-非守恒}
\end{align}
。在这种情形下，无法在分布论的框架内刻画弱解。然而，仍可通过在近似跳跃间断点处假定扮演 Rankine--Hugoniot 跳跃条件\eqref{align:ch08-rh} 角色的跳跃条件，在有界变差函数类 $BV$ 中引入弱解的概念。

合适的跳跃条件可由关于激波剖面的先验信息驱动：这些信息来自消失黏性方法、消失黏性--毛细论证，或相变理论中（所谓的）动力学关系。该理论的表述沿下述路线展开。

对\eqref{align:ch08-非守恒}，我们联系一个函数 $V(\tau;U_-,U_+)$，它定义在 $(-\infty,\infty)\times\mathbb R^n\times\mathbb R^n$ 上、取值于 $\mathbb R^n$，并具有下述性质：
\begin{align}
V(-\infty;U_-,U_+)=U_-,\qquad V(\infty;U_-,U_+)=U_+,
\label{align:ch08-边界剖面}\\
V(\tau;U,U)=U,
\label{align:ch08-对角剖面}\\
\big|\partial_\tau V(\tau;U_-,U_+)-\partial_\tau V(\tau;\bar U_-,\bar U_+)\big|
\leqslant a\big|(U_+-\bar U_+)-(U_--\bar U_-)\big|,
\label{align:ch08-Lipschitz剖面}
\end{align}
对所有 $U,U_-,U_+,\bar U_-,\bar U_+\in\mathbb R^n$、任一 $\tau\in(-\infty,\infty)$ 以及某个 $a>0$ 成立。进而要求：一个在左侧连接状态 $U_-$、在右侧连接状态 $U_+$、以速度 $s$ 传播的激波必须满足跳跃条件
\begin{align}
\int_{-\infty}^{\infty}A\big(V(\tau;U_-,U_+)\big)\partial_\tau V(\tau;U_-,U_+)\,\mathrm{d}\tau=s[U_+-U_-].
\label{align:ch08-非守恒跳跃}
\end{align}

在上述设定中，$V(\cdot;U_-,U_+)$ 表示激波剖面。注意在守恒情形 $A(U)=\mathrm{D}F(U)$ 下，不论 $V$ 如何选取，\eqref{align:ch08-非守恒跳跃} 都蜕化为 Rankine--Hugoniot 跳跃条件\eqref{align:ch08-rh}。

\hyperref[section:ch08-注记]{8.8 节}所引文献解释了上述装置如何在 $BV$ 函数框架内自然导出\eqref{align:ch08-非守恒} 的弱解概念。

\end{document}
